\documentclass{beamer}
\usepackage{hyperref}
\usepackage[russian]{babel}
\usepackage{fontspec}

\usepackage{latexsym,amsmath,xcolor,multicol,booktabs,calligra}
\usepackage{graphicx,listings,stackengine}
\usepackage{array}

\setmainfont{Liberation Serif}
\setsansfont{Liberation Sans}
\setmonofont{Liberation Mono}

\author{Камендов Максим, Морозов Даниил}
\title{Линковка в C++: эволюция, нагрузка, экосистема}
\subtitle{От Make до Cargo — как менялась сборка}
\institute{C++ Семинар}
\date{\today}
\usepackage{CNU}
\setbeamerfont{frame number}{family=\sffamily,size=\footnotesize}

\AtBeginSection[]{}
\AtBeginSubsection[]{}

\definecolor{deepblue}{rgb}{0,0,0.5}
\definecolor{deepred}{rgb}{0.6,0,0}
\definecolor{deepgreen}{rgb}{0,0.5,0}
\definecolor{halfgray}{gray}{0.55}
\definecolor{barred}{rgb}{0.9,0.3,0.2}
\definecolor{bargreen}{rgb}{0.2,0.6,0.2}

\lstset{
basicstyle=\ttfamily\small,
keywordstyle=\bfseries\color{deepblue},
emphstyle=\ttfamily\color{deepred},
stringstyle=\color{deepgreen},
breaklines=true,
columns=fullflexible,
numbers=left,
numberstyle=\small\color{halfgray},
rulesepcolor=\color{red!20!green!20!blue!20},
frame=shadowbox,
language=C++
}

\lstdefinelanguage{CMake}{
morekeywords={cmake_minimum_required,project,add_executable,target_link_libraries,add_library,target_include_directories,find_package,PUBLIC,PRIVATE,INTERFACE},
morecomment=[l]{\#},
morestring=[b]"
}

\lstdefinelanguage{Make}{
morecomment=[l]{\#},
morestring=[b]"
}

\lstdefinelanguage{TOML}{
morekeywords={package,dependencies,name,version,edition,build},
morecomment=[l]{\#},
morestring=[b]"
}

\newcommand{\progbar}[3]{%
\textcolor{#2}{\rule{#1}{4pt}}%
\hspace{2pt}%
\textcolor{halfgray}{\scriptsize #3}}

\begin{document}

\begin{frame}
\titlepage
\end{frame}

\begin{frame}{План}
\tableofcontents
\end{frame}

\section{Часть 1: Линковка}

\begin{frame}{Что на самом деле делает линкер}
\begin{itemize}
\item Линкер выполняет два принципиально разных прохода:
\begin{enumerate}
\item \textbf{Resolution} — сборка глобальной таблицы символов из всех объектных файлов.
Сильный (strong) символ — определение. Слабый (weak) — может быть переопределён.
\item \textbf{Relocation} — подстановка реальных адресов вместо временных меток
(секция + смещение).
\end{enumerate}
\item Современные линкеры: \texttt{ld} (GNU), \texttt{lld} (LLVM), \texttt{mold}.
\item Раздельная компиляция: каждый \texttt{.cpp} $\rightarrow$ \texttt{.o},
потом всё склеивается в один образ.
\end{itemize}
\end{frame}

\begin{frame}[fragile]{Анатомия объектного файла (ELF)}
\begin{itemize}
\item Основные секции:
\begin{itemize}
\item \texttt{.text} — исполняемый код
\item \texttt{.data} — инициализированные глобальные переменные
\item \texttt{.bss} — неинициализированные глобальные переменные
\item \texttt{.symtab} — таблица символов
\item \texttt{.rela.text} — таблица релокаций для секции кода
\end{itemize}
\item Strong vs Weak: в C++ глобальная переменная — strong, inline-функция — weak.
\item Диагностика:
\end{itemize}
\begin{lstlisting}
$ nm a.out | head -5
0000000000001129 T main
0000000000001144 W _Z3addii
                 U _ZNSolsEPFRSoS_E
\end{lstlisting}
\vspace{-0.3em}
\texttt{T} = strong defined, \texttt{W} = weak defined, \texttt{U} = undefined.
\end{frame}

\begin{frame}{Статическая и динамическая линковка}
\begin{itemize}
\item \textbf{Статическая} (\texttt{.a}): линкер копирует нужные секции прямо в
исполняемый файл. Больше размер, но нет runtime-зависимостей.
\item \textbf{Динамическая} (\texttt{.so}, \texttt{.dll}, \texttt{.dylib}):
линкер записывает только ссылку. Код загружается в память при запуске.
\item PIC (Position-Independent Code): обязателен для \texttt{.so}.
Позволяет загружать библиотеку по любому адресу без пересчёта релокаций.
\item Runtime linker (\texttt{ld.so}): \texttt{LD\_LIBRARY\_PATH},
\texttt{RPATH}, \texttt{RUNPATH}.
\end{itemize}
\end{frame}

\section{Часть 2: Линкер в C и C++}

\begin{frame}[fragile]{Name Mangling и ODR}
\begin{itemize}
\item C++ линкуется не как C: перегрузка функций, шаблоны, пространства имён
требуют \textbf{name mangling} — кодирования полного типа в имя символа.
\end{itemize}
\begin{lstlisting}
int add(int, int);         // _Z3addii
double add(double, double);// _Z3adddd
namespace ns {
  void foo();              // _ZN2ns3fooEv
}
\end{lstlisting}
\begin{itemize}
\item \texttt{extern "C"} отключает mangling — единственный способ
линковать C++ код с C-библиотеками.
\item ODR: каждое определение должно быть уникально, но для inline-функций,
шаблонов и \texttt{constexpr} сделано исключение — они порождают weak symbols.
\item Итог: простая программа на C++ может породить в десятки раз больше символов,
чем эквивалент на C.
\end{itemize}
\end{frame}

\begin{frame}[fragile]{Шаблоны, COMDAT и скрытый рост бинарника}
\begin{itemize}
\item Каждая инстанциация шаблона порождает weak symbol.
\item Линкер видит десятки копий \texttt{std::vector<int>::push\_back}
(по одной из каждого \texttt{.cpp}).
\item \textbf{COMDAT folding}: линкер удаляет дубликаты, оставляя одну копию.
Но сначала надо сгенерировать все копии, разложить их и сравнить.
\end{itemize}
\begin{lstlisting}
// a.cpp, b.cpp, c.cpp — каждая порождает:
void std::vector<int>::push_back(int&&);  // W sym
\end{lstlisting}
\begin{itemize}
\item Header-only библиотеки (Boost, nlohmann/json) эксплуатируют COMDAT.
Цена: рост времени компиляции и объёма .o файлов.
\item LTO (Link-Time Optimization) компенсирует: линкер перекомбинирует
код между единицами трансляции.
\end{itemize}
\end{frame}

\begin{frame}{C++20 Modules: попытка сломать архитектуру}
\begin{itemize}
\item Традиционная модель: \texttt{\#include} — текстовая вставка.
Один и тот же заголовок парсится N раз в N единицах трансляции.
\item Modules: \texttt{import} — уже не текст, а \textbf{BMI} (Binary Module Interface).
Препроцессор не нужен, AST готов.
\item С точки зрения линковки:
\begin{itemize}
\item Модуль компилируется один раз в BMI.
\item Импортёр получает уже готовый AST.
\item Меньше работы компилятору, но сборка становится жёстко
упорядоченной — нужно скомпилировать модуль до импортёра.
\end{itemize}
\item Статус на 2026: GCC, Clang и MSVC поддерживают, но миграция
существующих проектов идёт медленно. CMake уже умеет \texttt{target\_compile\_features(… MODULES)}.
\end{itemize}
\end{frame}

\section{Часть 3: Экосистема сборки}

\begin{frame}{От файлов к пакетам: три уровня абстракции}
\begin{center}
\small
\begin{tabular}{p{2cm}p{3.5cm}p{3.5cm}}
\toprule
\textbf{Уровень} & \textbf{Сущность} & \textbf{Инструмент} \\
\midrule
Файловый & Отдельные \texttt{.cpp}/\texttt{.h} файлы, timestamps & Make, Ninja \\
\midrule
Целевой & Target (библиотека/приложение) с транзитивными
свойствами & CMake, Meson, Bazel \\
\midrule
Пакетный & Пакет с версией, зависимостями и метаданными & Cargo, vcpkg, Conan \\
\bottomrule
\end{tabular}
\end{center}
\vspace{0.5em}
Ключевой вопрос: \textbf{что должен знать разработчик}, чтобы добавить зависимость?
\end{frame}

\begin{frame}[fragile]{Make: DAG над файловой системой}
\begin{itemize}
\item Концепция 1977 года (Stuart Feldman, Bell Labs).
\item Правило: «если цель старше зависимости — пересобрать».
\item Плюсы: языковая независимость, полный контроль.
\item Минусы: не знает про \texttt{\#include}, ручное описание графа,
никакой кроссплатформенности.
\end{itemize}
\begin{lstlisting}[language=Make,numbers=none]
app: main.o parser.o ; $(CXX) $^ -o $@
main.o: main.cpp parser.h ; $(CXX) -c main.cpp
# Добавили config.h в parser.h — нужно обновить строки!
\end{lstlisting}
Костыль: \texttt{gcc -MM} — внешний генератор зависимостей.
\end{frame}

\begin{frame}[fragile]{CMake: таргеты и транзитивные требования}
\begin{itemize}
\item CMake — генератор, а не сборочная система.
\item Главное изобретение: \textbf{Target} — абстракция «единица сборки»,
которая несёт с собой include-пути, флаги, зависимости.
\item PUBLIC / PRIVATE / INTERFACE — модель распространения свойств.
\item \texttt{find\_package}, \texttt{FetchContent} — встроенное управление
зависимостями.
\end{itemize}
\begin{lstlisting}[language=CMake]
add_library(network src/net.cpp include/net.h)
target_include_directories(network PUBLIC include/)

add_executable(server src/server.cpp)
target_link_libraries(server PUBLIC network)
# server автоматически получит -Iinclude
\end{lstlisting}
Цена: собственный язык с неочевидным синтаксисом.
\end{frame}

\begin{frame}[fragile]{Cargo: пакетный менеджер как сборочная система}
\begin{itemize}
\item Rust-экосистема: компилятор + линкер + пакетный менеджер = один бинарник.
\item Convention over Configuration: имя пакета = имя крейта.
\item Semver-резолвер: предотвращает конфликты транзитивных зависимостей.
\end{itemize}
\begin{lstlisting}[language=TOML,numbers=none]
[package]
name = "server"
version = "0.1.0"
edition = "2021"
[dependencies]
tokio = { version = "1", features = ["net"] }
\end{lstlisting}
\begin{itemize}
\item \texttt{cargo build} — одна команда на всё.
\item \texttt{build.rs} — кастомная логика на Rust вместо Make-шагов.
\item Ограничение: только Rust. Для C++ — комбинировать с CMake/vcpkg.
\end{itemize}
\end{frame}

\begin{frame}[fragile]{Одна задача — три подхода}
Задача: библиотека \texttt{greet} + приложение \texttt{hello}.

\small
\textbf{Make:}\qquad
\lstinline[language=Make]|hello: main.o libgreet.a ; $(CXX) $^ -o $@|

\textbf{CMake:}\qquad
\lstinline[language=CMake]|target_link_libraries(hello greet)|

\textbf{Cargo:}\qquad
\lstinline[language=TOML]|greet = { path = "../greet" }|

\normalsize
Эволюция: от явной команды \texttt{ar rcs} до декларативной строки.
Рост абстракции = уменьшение явного кода = меньше мест для ошибок.
\end{frame}

\section{Часть 4: Методология и результаты}

% \begin{frame}{Метрики сравнения}
% Методология: оцениваем каждый инструмент по 5 шкалам от 1 (плохо) до 5 (отлично).
%
% \medskip
% \begin{enumerate}
% \item \textbf{Абстракция} — насколько сущности инструмента соответствуют
% логике проекта (файл vs пакет).
% \item \textbf{Управление зависимостями} — как подключаются сторонние библиотеки.
% \item \textbf{Кроссплатформенность} — работает ли на Windows, Linux, macOS.
% \item \textbf{Диагностика} — насколько читаемы сообщения об ошибках.
% \item \textbf{Порог входа} — время до первой осмысленной сборки.
% \end{enumerate}
% \end{frame}

% \begin{frame}{Результаты сравнения}
% \centering
% \small
% \begin{tabular}{p{2cm}p{1.8cm}p{1.8cm}p{1.8cm}}
% \toprule
% \textbf{Метрика} & \textbf{Make} & \textbf{CMake} & \textbf{Cargo} \\
% \midrule
% Абстракция & 2 (файлы) & 4 (таргеты) & 5 (пакеты) \\
% \midrule
% Зависимости & 1 (ручные) & 3 (Fetch) & 5 (crates.io) \\
% \midrule
% Платформы & 2 (POSIX) & 5 (генератор) & 4 (встроено) \\
% \midrule
% Диагностика & 3 (честный Make) & 2 (свой язык) & 5 (Rust-сообщения) \\
% \midrule
% Порог входа & 4 & 2 & 5 \\
% \bottomrule
% \end{tabular}
%
% \medskip
% \textbf{Визуализация:}
% \bigskip
%
% Make:~\progbar{2cm}{barred}{файлы}
%
% CMake:~\progbar{3.2cm}{deepblue}{таргеты}
%
% Cargo:~\progbar{4.5cm}{bargreen}{пакеты}
%
% \hfill \scriptsize Длина = (балл по шкале абстракции) $\times$ 1cm
% \end{frame}

\begin{frame}{Выводы}
\begin{itemize}
\item \textbf{Make}: DAG-модель пересборки, но без знаний о C++ —
никаких транзитивных зависимостей или типов целей.
\item \textbf{CMake}: таргеты с транзитивными свойствами —
лучший компромисс для C++ сегодня. Минус: сложный синтаксис.
\item \textbf{Cargo}: менеджер пакетов + компилятор + линкер
в одном бинарнике. Минус: привязан к одному языку.
% \item C++ платит исторический долг: \texttt{\#include}, ODR, нет
% модульности. Modules (C++20) + vcpkg/Conan сокращают разрыв,
% но инерция архитектуры огромна.
% \item Правило: выбирай наивысший уровень абстракции, доступный
% для твоего языка и инфраструктуры.
\end{itemize}
\end{frame}

\section{Источники}

\begin{frame}[allowframebreaks]{Источники}
\nocite{iso_cpp_odr,cpp20_modules,gnu_make,cmake_docs,cargo_book,elf_spec,itanium_abi,mold_docs,gcc_lto}
\bibliographystyle{unsrt}
\bibliography{ref}
\end{frame}

\end{document}
